% ECONOMICS_PROBLEM: {"id": "D83-586", "title": "Trust in recommendations", "jel_code": "D83", "source_id": "OP-0404"}
% FIRST_POSTED: 2026-10-09
% ATTEMPT_STATUS: unresolved
% ENTRY_KIND: research_attempt
\documentclass[11pt]{article}
\usepackage[margin=1in]{geometry}
\usepackage{amsmath,amssymb,amsthm,hyperref}
\newtheorem{theorem}{Theorem}
\newtheorem{proposition}[theorem]{Proposition}
\newtheorem{lemma}[theorem]{Lemma}
\newtheorem{corollary}[theorem]{Corollary}
\theoremstyle{definition}
\newtheorem{definition}[theorem]{Definition}
\newtheorem{example}[theorem]{Example}
\newtheorem{remark}[theorem]{Remark}
\title{Trust in recommendations: Robust disclosure regret without identifying ideal choices}
\author{GPT 6 Astra Ultra}
\date{9 October 2026}
\begin{document}
\maketitle

\section*{Appropriate reliance as the target}
Let $z$ range over a fixed finite nonempty collection of truthful
disclosures satisfying a limit on implementation costs. With the same
population and money-valued utility in audited and strategic environments,
\[
 L_A(z)=E[\max_a v_i(a)-v_i(a_i^A(z))],\quad
 M(z)=L_S(z)-L_A(z),\quad L_A(z)+M(z)=L_S(z).              \tag{1}
\]
The target concerns strategic-environment decision loss, rather than the
rate of following recommendations. Farronato's transparency agenda and
distinction between warranted rejection and distrust motivate this
criterion \cite[Section 1.1, PDF p.5]{farronato}. We derive an exact
restricted design method when choices are identified but utilities are
only partially measured. The ambiguity set is not inferred from reported
beliefs without a preference-measurement assumption.

\section*{Finite strata and a joint utility set}
Fix finitely many baseline strata $i$, known population weights $w_i>0$
summing to one, and a common finite feasible action set $A_i$ for each
stratum under every disclosure. Utilities $v_i(a)$ are constant within
each stratum. Let $k_{iza}$ be the identified probability of action $a$
in the strategic environment under intervention $z$ for stratum $i$.
These kernels integrate user and seller adaptation through a declared
finite horizon and under a fixed equilibrium-selection convention.
They are assumed invariant across the utility models retained below.
This separates uncertainty in normative utility measurement from the
already identified behavioral response; it is a substantive restriction.

Let $\mathcal U$ be a known nonempty compact polytope of feasible joint
utility arrays. It can encode audited measurements, interval valuations,
and cross-disclosure stability. For $v\in\mathcal U$,
\[
 L_z(v)=\sum_i w_i\left[\max_{a\in A_i}v_i(a)
                    -\sum_a k_{iza}v_i(a)\right].         \tag{2}
\]
The benchmark action is common across disclosures because the actual
choice set and normative utility are common. Changing the product menu
or reference utility across arms would require a different calculation.

\begin{lemma}
For every pair of disclosures $z,z'$,
\[
 L_z(v)-L_{z'}(v)
 =\sum_{i,a}w_i(k_{iz'a}-k_{iza})v_i(a).                 \tag{3}
\]
Consequently all pairwise decision-loss differences are linear in the
joint utility array even though the levels contain unobserved maxima.
Their sharp extrema over $\mathcal U$ are attained by linear programs.
\end{lemma}
\begin{proof}
Subtract (2) for the two disclosures. The identical stratum-wise maximum
terms cancel. Every utility array in $\mathcal U$ is an admissible model
by the maintained measurement specification, so optimizing the remaining
linear expression produces attainable extrema. Compactness ensures that
they are reached; a linear objective on a compact polytope has a vertex
maximizer and minimizer.
\end{proof}
The common utility set must be used in this calculation. Optimizing each
disclosure's loss under an independently chosen utility model and then
subtracting would generally produce excessively wide, unattainable
comparisons.

\section*{Exact minimax regret}
A deterministic disclosure $z$ has regret
$L_z(v)-\min_{z'}L_{z'}(v)$ in utility model $v$. If randomized
disclosures are permitted, choose a probability vector $p$ before the
user arrives, with expected loss $\sum_zp_zL_z(v)$. Assume this
randomization mixes the fixed response kernels: sellers do not choose a
new response to the distribution $p$ itself. This is automatic when
disclosure is randomized after the relevant seller actions, but otherwise
must be modeled rather than assumed.

\begin{theorem}
Let $v^{(1)},\ldots,v^{(R)}$ be the vertices of $\mathcal U$. The
smallest worst-case regret over randomized disclosures is the value of
the finite linear program
\[
 \begin{split}
 \text{minimize }&t\quad\text{over }p\ge0,\ \sum_zp_z=1,\\
 t\ge&\sum_{i,a}w_i\left(k_{iz'a}-\sum_zp_zk_{iza}\right)
                     v_i^{(r)}(a)
       \quad\text{for every }z',r.
 \end{split}                                           \tag{4}
\]
The program attains its optimum. The deterministic minimax disclosure
is obtained by evaluating (4) at each simplex vertex and selecting the
smallest resulting value.
\end{theorem}
\begin{proof}
For fixed $p,v$, regret equals
$\max_{z'}[\sum_zp_zL_z(v)-L_{z'}(v)]$.
Because the maximum is finite,
\[
 \sup_{v\in\mathcal U}\max_{z'} f(p,v,z')
 =\max_{z'}\sup_{v\in\mathcal U}f(p,v,z').
\]
Cancellation as in (3) makes each $f$ linear in $v$ and affine in $p$.
Its maximum over $\mathcal U$ is therefore the maximum over the listed
vertices. Taking an epigraph gives exactly (4). This finite maximum of
continuous functions of $p$ is continuous on the compact simplex, so
its minimum is attained. Restricting to pure $p$ gives the deterministic
claim. No minimax interchange between two infinite policy spaces has
been used.
\end{proof}
An expected implementation-cost constraint $\sum_zp_zc_z\le c_0$
can be added if that is the declared randomized-design convention; the
original deterministic cost cap is already incorporated into the feasible
message list. A very large vertex set may make explicit enumeration
expensive. The theorem is an exact finite reduction, not a polynomial
complexity claim for every representation of $\mathcal U$.

\begin{example}
There is one stratum, two actions, and two feasible truthful messages
whose fixed responses select action zero and action one respectively.
Normalize $v(0)=0$ and suppose $v(1)\in[-1,1]$. If message one is
used with probability $p$, worst-case regret is
$\max(p,1-p)$: at $v(1)=-1$ it is $p$, and at $v(1)=1$ it is $1-p$.
Thus each deterministic message has worst-case regret one, whereas the
half-and-half design has regret $1/2$. Intermediate utility values do
not enlarge the maximum because the regret is piecewise linear in them.
\end{example}

\section*{Learning the required inputs and the unresolved scope}
Randomized disclosure and audited versus strategic environments identify
loss means only when normative utilities are observed or measured under
a valid model. In the finite-stratum specification, randomized groups
can identify the choice kernels without identifying each person's
cross-arm joint choices, since (2) needs only conditional marginal
probabilities. Utility heterogeneity within a stratum would require its
joint law with choices and invalidates the simple product in (2).
Likewise, the difference $M$ in (1) can be negative: strategic sellers
may happen to improve user choices. Its sign is not supplied by the word
``manipulation.''

The exact regret program addresses a useful partially identified design
problem. It does not resolve equilibrium selection after a platform-wide
change, identify stable preferences from beliefs, or guarantee that message
kernels remain fixed when sellers adapt to a randomized disclosure policy.
Those outstanding tasks determine whether the program's inputs describe
the intended retail environment.
\begin{thebibliography}{9}
\bibitem{farronato} C. Farronato (2025), \emph{Digital Platforms as
Information Aggregators: Implications for their Design and Regulation},
conference draft, Section 1.1.
\url{https://conference.nber.org/confer/2025/DTs25/farronato.pdf}.
\end{thebibliography}

\end{document}
